% ============================================================================
%  Chapitre 2 — Contexte industriel et analyse de l'existant
%  Cible : 8 à 10 pages. À rédiger J2-J3.
% ============================================================================
\chapter{Contexte industriel et analyse de l'existant}
\label{chap:contexte}

\section{L'entreprise et l'entité d'accueil}
\label{sec:ctx-entreprise}

\aecrire{1,5 page maximum. Ne pas faire une plaquette commerciale : le jury
attend le \emph{contexte organisationnel du projet}, pas le chiffre
d'affaires. Points à couvrir : activité de la division, positionnement du
produit dans la gamme, composition et localisation des équipes,
articulation entre équipe de développement firmware et équipe DevOps,
rattachement de l'auteur.}

\section{Le produit et son logiciel embarqué}
\label{sec:ctx-produit}

\subsection{Fonction du produit}
\aecrire{Rôle d'une centrale de mesure / compteur d'énergie communicant :
acquisition des grandeurs électriques, journalisation, exposition des données
aux systèmes de supervision.}

\subsection{Architecture matérielle}
La cible matérielle repose sur un \gls{soc} Texas~Instruments AM642 intégré sur
un \gls{som}, associé à une carte porteuse. La configuration de construction
désigne cette cible par la machine Yocto \texttt{\Carte}.

\aecrire{Décrire : cœurs Cortex-A53 (Linux) et cœurs Cortex-R5F (temps réel),
mémoire, stockage eMMC, présence d'un \gls{tpm}. Insister uniquement sur ce
qui \emph{contraint la chaîne de construction} : compilation croisée
aarch64, firmware R5F à générer avec une chaîne d'outils propriétaire,
signature et chaîne de confiance au démarrage.}

\subsection{Pile logicielle}
\label{ssec:ctx-pile-logicielle}

La pile logicielle se décompose en trois couches, chacune associée à un rythme
de construction et à des outils différents :

\begin{table}[H]
\centering
\caption{Décomposition de la pile logicielle embarquée}
\label{tab:ctx-pile}
\begin{tabularx}{\textwidth}{@{}l X l@{}}
\toprule
\textbf{Couche} & \textbf{Contenu} & \textbf{Construction} \\
\midrule
Distribution & Noyau Linux, \gls{bsp}, paquets système, image disque &
  Yocto / kas \\
Sécurité & Brique de cybersécurité transverse (\gls{csb}) & Meson \\
Application & Services métier \Firmware{} (gestion d'équipement, de défauts,
  interfaces nord) & CMake \\
\bottomrule
\end{tabularx}
\end{table}

\aecrire{Décrire brièvement les services applicatifs : gestion
administrative et système, gestion des équipements, gestion des défauts,
notifications, et les interfaces nord Modbus, BACnet et \gls{edm}. Mentionner
l'organisation en dépôts multiples avec sous-modules Git, car c'est une
contrainte forte pour la chaîne d'intégration.}

\section{Organisation des dépôts}
\label{sec:ctx-depots}

\aecrire{Expliquer la séparation en dépôts : configuration d'image,
couches BSP et distribution, application, tests d'automatisation, images de
construction, configuration des exécuteurs, supervision. Justifier pourquoi
cette séparation complique l'intégration (versions croisées, cohérence des
références, secrets d'accès inter-dépôts). Un tableau des dépôts est plus
lisible qu'un paragraphe.}

\begin{table}[H]
\centering
\caption{Dépôts intervenant dans la chaîne de construction}
\label{tab:ctx-depots}
\begin{tabularx}{\textwidth}{@{}>{\raggedright\arraybackslash}p{4cm} X l@{}}
\toprule
\textbf{Dépôt} & \textbf{Rôle} & \textbf{Branche} \\
\midrule
\texttt{ci} & Workflows et scripts d'intégration & \texttt{main} \\
\texttt{docker-image} & Image de l'environnement de compilation & \texttt{main} \\
\texttt{evo-image-config} & Configuration kas / Yocto de la distribution & \texttt{scarthgap} \\
\texttt{sdp} & Code applicatif et sous-modules & \texttt{develop} \\
\texttt{automation-tests} & Suites de tests d'intégration & \texttt{main} \\
\texttt{runner-config} & Infrastructure des exécuteurs & \texttt{Prod} \\
\texttt{monitoring-ci} & Exportateurs de métriques & \texttt{main} \\
\bottomrule
\end{tabularx}
\end{table}

\section{Analyse de l'existant}
\label{sec:ctx-existant}

\aecrire{C'est LA section qui justifie tout le mémoire. Sans un « avant »
décrit honnêtement, le chapitre~6 n'a pas de point de comparaison. Décrire
la situation initiale sans la caricaturer.}

\subsection{Pratiques de construction initiales}
\aecrire{Construction sur poste de développeur, environnements divergents
d'un poste à l'autre, dépendances installées manuellement, absence de
version de référence de la chaîne d'outils.}

\subsection{Pratiques de test et de qualité initiales}
\aecrire{Tests unitaires exécutés à la demande, absence de mesure de
couverture historisée, analyse statique lancée ponctuellement en fin de cycle,
analyse de composants réalisée manuellement avant les jalons.}

\subsection{Limites identifiées}
\label{ssec:ctx-limites}

\aecrire{Formuler chaque limite de façon à ce qu'elle soit reprise telle
quelle comme exigence au chapitre~4 et comme critère d'évaluation au
chapitre~6. Une limite = une exigence = une mesure. Conserver cette
traçabilité, le jury la cherchera.}

\begin{table}[H]
\centering
\caption{Limites de la situation initiale et exigences induites}
\label{tab:ctx-limites}
\begin{tabularx}{\textwidth}{@{}l X l@{}}
\toprule
\textbf{Réf.} & \textbf{Limite constatée} & \textbf{Exigence} \\
\midrule
L1 & Environnements de construction divergents entre développeurs & EF1 \\
L2 & Détection tardive des régressions & EF2 \\
L3 & Analyse statique non systématique & EF4 \\
L4 & Inventaire des composants tiers non tenu à jour & EF5 \\
L5 & Absence de mesure historisée de la qualité & EF6 \\
L6 & Administration manuelle des machines d'exécution & ENF3 \\
\bottomrule
\end{tabularx}
\end{table}

\section{Synthèse}
\label{sec:ctx-synthese}

\aecrire{Un paragraphe de transition vers le chapitre~3 : avant de concevoir,
examiner ce que la littérature et l'industrie proposent pour ces problèmes.}
